\originalchapter{习题课九: 受迫RLC电路}
题源为\href{https://ocw.mit.edu/courses/8-02-physics-ii-electricity-and-magnetism-spring-2007/resources/solving09/}{Problem Solving 9}, PDF第1--10页(原页脚沿用Solving 10). Q1--Q15及考试练习题(a)--(e)全部保留. 约定$V(t)=V_0\sin\omega t$, $I(t)=I_0\sin(\omega t-\phi)$, $\omega>0$. 所有解答为AI独立解答.

\begin{exercise}
纯电阻电路: Q1求电流幅度、相位; Q2在一般串联RLC幅度公式中取哪些$L,C$能回到本题? Q3求平均耗散功率.
\end{exercise}
\begin{solution}
Q1. $I=V/R$, $I_0=V_0/R$, $\phi=0$.
Q2. $L=0,C\to\infty$, 使两个电抗为零; $C=0$会断路, 不能用来表示没有串联电容.
Q3. $\langle P\rangle=V_0^2\langle\sin^2\omega t\rangle/R=V_0^2/(2R)$.
\end{solution}

\begin{exercise}
纯电感电路: Q4解$L\dot I=V(t)$, 给正弦电流幅度和相位; Q5在一般公式中选哪些$R,C$? Q6求周期平均功率.
\end{exercise}
\begin{solution}
Q4. $I=-(V_0/\omega L)\cos\omega t+I_{\rm dc}$, 纯正弦项幅度$V_0/(\omega L)$、$\phi=\pi/2$. 理想无电阻电感还允许任意常数初始电流, 它不会自行衰减; 指定零直流分量才得到原题所用纯正弦形式.
Q5. $R=0,C\to\infty$.
Q6. 正弦项功率正负交替, 平均$\sin\omega t\cos\omega t=0$, 常数项乘电压的平均也为零, 所以$\langle P\rangle=0$. 不代表$U_B$一直不变.
\end{solution}

\begin{exercise}
纯电容电路: Q7由$q/C=V(t)$求电流幅度、相位; Q8求平均功率.
\end{exercise}
\begin{solution}
Q7. $I=\dot q=\omega CV_0\cos\omega t$, $I_0=\omega CV_0$, $\phi=-\pi/2$, 电流领先电压.
Q8. $P=\omega CV_0^2\sin\omega t\cos\omega t$, 周期平均为零, 电容交替吸收和返还电场能.
\end{solution}

\begin{exercise}
黑箱一. 电源、电阻$R$和仅含一个电感或一个电容的黑箱串联. 驱动幅度$V_0=100\sqrt2\,\mathrm V$, $\omega=10\,\mathrm{rad/s}$, 实测$I(t)=(10\,\mathrm A)\sin(\omega t+\pi/4)$.
Q9电流领先还是落后? Q10箱中是什么? Q11求$R$数值与单位. Q12求电感或电容值与单位.
\end{exercise}
\begin{solution}
Q9. 领先$\pi/4$. Q10. 对应$\phi=-\pi/4$, 电抗为负, 所以是电容.
Q11. $Z=V_0/I_0=10\sqrt2\,\Omega$, $R=Z\cos\phi=10\,\Omega$.
Q12. 有符号电容电抗$X_C=-1/(\omega C)=-R=-10\,\Omega$, 因而$1/(\omega C)=10\,\Omega$, $C=0.01\,\mathrm F$. 原图串联结构是这一判断的条件, 不能把结果套给并联黑箱.
\end{solution}

\begin{exercise}
黑箱二. 电源、电阻$R=3\,\Omega$和含电感、电容之一或两者的黑箱串联. $V_0=6\,\mathrm V$, 在$\omega_1=1\,\mathrm{rad/s}$时电流与电压同相, 在$\omega_2=2\,\mathrm{rad/s}$时相差$\pi/4$.
Q13箱中有什么, 第二频率时电流领先还是落后? Q14求元件数值和单位. Q15求第一频率的平均耗散功率.
\end{exercise}
\begin{solution}
Q13. 有限正$L$单独给正电抗, 有限正$C$单独给负电抗, 都不能在有限非零频率抵消. 故两者都有, 且$\omega_1L=1/(\omega_1C)$. 第二频率更高, 电感项占优, 电流落后.
Q14. $\omega_2L-1/(\omega_2C)=R\tan(\pi/4)=3\,\Omega$. 用$1/C=\omega_1^2L$消去, $(2-1/2)L=3$按相应SI数值求得$L=2\,\mathrm H$, $C=0.5\,\mathrm F$.
Q15. 第一频率$Z=R$, $I_0=2\,\mathrm A$, 平均功率$I_0^2R/2=6\,\mathrm W$.
\end{solution}

\begin{exercise}
考试练习题. 电池$V$、开关和电阻$R$串联后到一个节点, 节点分两路: 中路电容$C$, 右路电阻$R$与电感$L$串联, 两路返回电池负端. $i_1$经左电阻向右, $i_2$经电容向下, $i_3$经右电阻和电感向下. 初始电容未充电且所有电流零, 合闸.
(a)外圈(含电池、电感)顺时针写电势变化总和.
(b)刚合闸求三个电流, 左圈电感为零.
(c)很久后求三个电流.
(d)很久后求电容电荷.
(e)再次断开开关后瞬间求三个电流, 左圈电感仍为零.
\end{exercise}
\begin{solution}
(a)$V-Ri_1-Ri_3-L\dot i_3=0$.
(b)电容电压初值零、电感电流初值零, 所以$i_1=V/R,i_2=V/R,i_3=0$.
(c)直流稳态电容开路、电感短路, $i_1=i_3=V/(2R),i_2=0$.
(d)节点电压$V/2$, 所以$q=CV/2$.
(e)电池支路断开,$i_1=0$, 电感电流不能跳变, $i_3=V/(2R)$. 节点电荷守恒给$i_2=-i_3=-V/(2R)$, 即电容开始供给右支路, 而不是电感电流突然消失.
\end{solution}
